\renewcommand{\currentchapter}{5}
\title{Evaluasi Program UTBK 2026}
\subtitle{Bab 5 \textbar\ SMAN 1 Muntilan: capaian, kesenjangan, dan pelajaran}
\date{SMAN 1 Muntilan bersama StudentPro \textbar\ Tahun Ajaran 2026/2027}

\begingroup
\setbeamercolor{background canvas}{bg=StudentNavy}
\begin{frame}[plain]
  \vspace{0.65cm}
  \kicker{SMAN 1 Muntilan \textbar\ StudentPro}
  \vspace{0.08cm}
  {\usebeamerfont{title}\color{StudentWhite}\inserttitle\par}
  \vspace{0.38cm}
  {\usebeamerfont{subtitle}\color{StudentSoft}\insertsubtitle\par}
  \vfill
  {\color{StudentGold}\rule{2.3cm}{2.5pt}}\par
  \vspace{0.20cm}
  {\usebeamerfont{author}\color{StudentWhite}\insertauthor}\par
  {\usebeamerfont{date}\color{StudentSoft}\insertdate}
  \vspace{0.42cm}
\end{frame}
\endgroup

\begin{frame}{Evaluasi menjawab lebih dari pertanyaan lulus atau tidak}
  \kicker{Tujuan Bab 5}
  \begin{columns}[T,onlytextwidth]
    \begin{column}{0.30\textwidth}
      \begin{block}{Capaian}
        Apa yang telah dihasilkan oleh rangkaian pembimbingan 2026?
      \end{block}
    \end{column}
    \begin{column}{0.30\textwidth}
      \begin{block}{Kesenjangan}
        Data, proses, dan hasil apa yang belum sesuai kebutuhan program?
      \end{block}
    \end{column}
    \begin{column}{0.30\textwidth}
      \begin{block}{Pelajaran}
        Perubahan apa yang harus dilakukan untuk pembimbingan 2027?
      \end{block}
    \end{column}
  \end{columns}
  \vspace{0.28cm}
  \claim{Evaluasi yang baik menghasilkan keputusan, bukan sekadar daftar keberhasilan.}
\end{frame}

\begin{frame}[shrink=10]{Empat lapisan digunakan untuk membaca program}
  \kicker{Kerangka Evaluasi}
  {\fontsize{10.3}{12.8}\selectfont
  \renewcommand{\arraystretch}{1.00}
  \begin{tabularx}{\textwidth}{>{\bfseries\color{StudentBlue}}p{0.18\textwidth} >{\RaggedRight\arraybackslash}p{0.31\textwidth} >{\RaggedRight\arraybackslash}X}
    \toprule
    \textbf{Lapisan} & \textbf{Pertanyaan} & \textbf{Contoh bukti} \\
    \midrule
    Input & Apa yang disiapkan? & Materi, bank soal, tentor, LMS, dan jadwal \\
    Proses & Apa yang benar-benar dilakukan? & Latihan, tryout, pembahasan, konsultasi, monitoring \\
    Output & Apa hasil langsungnya? & Ketuntasan, skor, tren, rencana tindak lanjut \\
    Outcome & Apa dampak akhirnya? & Penerimaan, ketepatan pilihan, dan kesiapan studi \\
    \bottomrule
  \end{tabularx}}
  \vspace{0.15cm}
  \centering{\color{StudentNavy}\bfseries\fontsize{10.8}{13.5}\selectfont Outcome tidak dapat dijelaskan tanpa membaca kualitas input, proses, dan output.}
\end{frame}

\begin{frame}{Kekuatan utama program terletak pada kelengkapan siklus}
  \kicker{Temuan Positif}
  \begin{columns}[T,onlytextwidth]
    \begin{column}{0.46\textwidth}
      \begin{itemize}
        \item Program tidak berhenti pada penyampaian materi.
        \item Bank soal dan tryout memberi kesempatan mengukur kemampuan.
        \item Pembahasan dan remedial menindaklanjuti kesalahan.
      \end{itemize}
    \end{column}
    \begin{column}{0.48\textwidth}
      \begin{itemize}
        \item Konsultasi menghubungkan hasil belajar dengan pilihan studi.
        \item Monitoring menjaga kegiatan di antara pertemuan.
        \item StudentPro menyediakan dasar integrasi data program.
      \end{itemize}
    \end{column}
  \end{columns}
  \vspace{0.22cm}
  \begin{beamercolorbox}[rounded=true,sep=0.20cm,wd=\textwidth]{block body}
    \centering\bfseries Fondasi program 2026 sudah membentuk alur belajar, mengukur, memperbaiki, dan mendampingi.
  \end{beamercolorbox}
\end{frame}

\begin{frame}{Konsultasi menunjukkan jangkauan pendampingan individual}
  \kicker{Capaian Proses yang Terdokumentasi}
  \begin{columns}[T,onlytextwidth]
    \begin{column}{0.30\textwidth}
      \metric{124}{siswa tercatat menerima konsultasi akademik}
    \end{column}
    \begin{column}{0.64\textwidth}
      \begin{itemize}
        \item Diagnostik dan strategi belajar.
        \item Evaluasi hasil tryout.
        \item Pemilihan jurusan dan kampus.
        \item Rekap hasil serta tindak lanjut.
      \end{itemize}
      \vspace{0.12cm}
      {\color{StudentGreen}\bfseries\fontsize{10.5}{13}\selectfont Kekuatan: layanan tidak hanya bersifat klasikal.}
    \end{column}
  \end{columns}
  \vspace{0.18cm}
  {\color{StudentGray}\fontsize{9.5}{11.8}\selectfont Angka 124 adalah jumlah siswa penerima layanan yang terdokumentasi, bukan jumlah sesi dan bukan jumlah seluruh peserta program.}
\end{frame}

\begin{frame}{Rekap yang tersedia menunjukkan capaian lintas bidang}
  \kicker{Hasil Seleksi yang Terdokumentasi}
  \begin{columns}[T,onlytextwidth]
    \begin{column}{0.36\textwidth}
      \metric{8}{catatan penerimaan melalui jalur SNBT 2026}
      \vspace{0.12cm}
      {\color{StudentGray}\fontsize{9.5}{11.8}\selectfont Data menggunakan rekap internal yang tersedia saat laporan disusun.}
    \end{column}
    \begin{column}{0.58\textwidth}
      \begin{itemize}
        \item Hukum dan Manajemen.
        \item Kedokteran dan Teknologi Pangan.
        \item Sains Data.
        \item Teknik Sipil: dua catatan.
        \item Geologi.
      \end{itemize}
    \end{column}
  \end{columns}
  \vspace{0.18cm}
  \begin{beamercolorbox}[rounded=true,sep=0.18cm,wd=\textwidth]{block body}
    \centering\bfseries Capaian mencakup saintek, soshum, kesehatan, data, dan rekayasa.
  \end{beamercolorbox}
\end{frame}

\begin{frame}{Jalur mandiri melengkapi gambaran capaian}
  \kicker{Rekap UM UGM 2026}
  \begin{columns}[T,onlytextwidth]
    \begin{column}{0.34\textwidth}
      \metric{5}{catatan penerimaan melalui UM UGM}
    \end{column}
    \begin{column}{0.60\textwidth}
      \begin{itemize}
        \item Program vokasi bidang pembangunan dan sipil.
        \item Teknik Fisika.
        \item Rekayasa infrastruktur.
        \item Pariwisata.
      \end{itemize}
      \vspace{0.12cm}
      {\color{StudentBlue}\bfseries\fontsize{10.5}{13}\selectfont Makna: pembimbingan pilihan studi tidak berhenti setelah SNBT.}
    \end{column}
  \end{columns}
  \vspace{0.20cm}
  {\color{StudentGray}\fontsize{9.5}{11.8}\selectfont Rekap jalur mandiri dilaporkan terpisah agar tidak dianggap sebagai hasil SNBT.}
\end{frame}

\begin{frame}{Posisi 53 harus disertai ruang lingkup pemeringkatan}
  \kicker{Catatan tentang Ranking}
  \begin{columns}[T,onlytextwidth]
    \begin{column}{0.30\textwidth}
      \metric{53}{posisi yang tercantum pada catatan internal UTBK 2026}
    \end{column}
    \begin{column}{0.64\textwidth}
      \begin{block}{Yang masih harus dilengkapi}
        Wilayah pembanding, populasi sekolah, indikator pembentuk ranking, tanggal data, dan lembaga yang menerbitkan.
      \end{block}
      \vspace{0.12cm}
      {\color{StudentGold}\bfseries\fontsize{10.5}{13}\selectfont Perlakuan laporan: dicatat sebagai data internal, belum digunakan sebagai klaim publik tanpa metadata.}
    \end{column}
  \end{columns}
\end{frame}

\begin{frame}{Target 70 persen belum boleh dibandingkan secara langsung}
  \kicker{Masalah Penyebut}
  \begin{columns}[T,onlytextwidth]
    \begin{column}{0.42\textwidth}
      \claim{Persentase keberhasilan memerlukan jumlah peserta sasaran yang final.}
      \vspace{0.18cm}
      {\color{StudentGray}\fontsize{10.5}{13}\selectfont Catatan penerimaan tidak otomatis sama dengan seluruh peserta pembimbingan.}
    \end{column}
    \begin{column}{0.52\textwidth}
      \begin{block}{Rumus yang benar}
        \[
          \text{Capaian}=
          \frac{\text{peserta diterima (unik)}}
          {\text{kohor sasaran (unik)}}\times 100\%
        \]
      \end{block}
      \vspace{0.10cm}
      {\color{StudentGold}\bfseries\fontsize{10.3}{12.8}\selectfont Status: menunggu daftar kohor dan hasil final yang tervalidasi.}
    \end{column}
  \end{columns}
\end{frame}

\begin{frame}{Data capaian perlu dibedakan berdasarkan jalur}
  \kicker{Mencegah Penghitungan Ganda}
  {\fontsize{10.2}{12.7}\selectfont
  \renewcommand{\arraystretch}{1.00}
  \begin{tabularx}{\textwidth}{>{\bfseries\color{StudentBlue}}p{0.22\textwidth} >{\RaggedRight\arraybackslash}X >{\RaggedRight\arraybackslash}p{0.28\textwidth}}
    \toprule
    \textbf{Kelompok} & \textbf{Cara mencatat} & \textbf{Risiko bila dicampur} \\
    \midrule
    SNBT & Hasil seleksi nasional berdasarkan UTBK & Dianggap sama dengan jalur mandiri \\
    Jalur mandiri & Dicatat per PTN dan tahap seleksi & Keberhasilan UTBK dapat dinilai berlebihan \\
    Diterima ganda & Gunakan identitas peserta unik & Satu siswa dihitung sebagai beberapa keberhasilan \\
    Belum ada hasil & Tetap masuk denominator bila bagian kohor & Persentase menjadi bias \\
    \bottomrule
  \end{tabularx}}
\end{frame}

\begin{frame}{Kesenjangan terbesar berada pada integrasi data}
  \kicker{Temuan yang Perlu Diperbaiki}
  \begin{columns}[T,onlytextwidth]
    \begin{column}{0.46\textwidth}
      \begin{itemize}
        \item Daftar kohor peserta belum menjadi satu sumber data final.
        \item Jumlah sesi belum terpisah jelas dari jumlah siswa.
        \item Data tryout, konsultasi, dan hasil seleksi masih tersebar.
      \end{itemize}
    \end{column}
    \begin{column}{0.48\textwidth}
      \begin{itemize}
        \item Metadata ranking belum lengkap.
        \item Tren per subtes belum seragam untuk seluruh siswa.
        \item Tindak lanjut belum selalu memiliki status selesai dan bukti ulang.
      \end{itemize}
    \end{column}
  \end{columns}
  \vspace{0.22cm}
  \claim{Program sudah kaya kegiatan; evaluasinya perlu satu bahasa data.}
\end{frame}

\begin{frame}[shrink=5]{Kualitas proses harus terlihat sebelum hasil seleksi keluar}
  \kicker{Indikator Evaluasi Berkala}
  {\fontsize{9.9}{12.3}\selectfont
  \renewcommand{\arraystretch}{0.97}
  \begin{tabularx}{\textwidth}{>{\bfseries\color{StudentBlue}}p{0.21\textwidth} >{\RaggedRight\arraybackslash}X >{\RaggedRight\arraybackslash}p{0.30\textwidth}}
    \toprule
    \textbf{Indikator} & \textbf{Ukuran} & \textbf{Makna} \\
    \midrule
    Cakupan diagnostik & Persentase kohor yang memiliki baseline & Tidak ada siswa memulai tanpa peta awal \\
    Ketuntasan belajar & Materi dan latihan selesai tepat waktu & Konsistensi proses \\
    Kemajuan & Perubahan akurasi dan waktu per subtes & Perbaikan yang dapat dijelaskan \\
    Tindak lanjut & Remedial selesai dan diuji ulang & Data menghasilkan tindakan \\
    Konsultasi & Masalah, keputusan, dan status tindak lanjut & Pendampingan individual dapat ditelusuri \\
    \bottomrule
  \end{tabularx}}
\end{frame}

\begin{frame}[shrink=7]{StudentPro perlu menjadi sumber data utama program}
  \kicker{Perbaikan Sistem 2027}
  \begin{columns}[T,onlytextwidth]
    \begin{column}{0.46\textwidth}
      \begin{itemize}
        \item Tetapkan kohor UTBK 2027 dan identitas peserta unik.
        \item Simpan baseline serta target per subtes.
        \item Beri label materi, indikator, dan konteks LBI.
        \item Rekam waktu, akurasi, serta jenis kesalahan.
      \end{itemize}
    \end{column}
    \begin{column}{0.48\textwidth}
      \begin{itemize}
        \item Hubungkan remedial dengan tes ulang.
        \item Catat konsultasi dan status tindak lanjut.
        \item Verifikasi hasil seleksi per jalur.
        \item Sediakan laporan siswa, kelas, dan orang tua.
      \end{itemize}
    \end{column}
  \end{columns}
  \vspace{0.16cm}
  \begin{beamercolorbox}[rounded=true,sep=0.18cm,wd=\textwidth]{block body}
    \centering\bfseries Satu siswa, satu riwayat, banyak bukti perkembangan.
  \end{beamercolorbox}
\end{frame}

\begin{frame}{Laporan bulanan harus berujung pada keputusan}
  \kicker{Siklus Evaluasi 2027}
  {\fontsize{10.0}{12.5}\selectfont
  \renewcommand{\arraystretch}{0.99}
  \begin{tabularx}{\textwidth}{>{\bfseries\color{StudentBlue}}p{0.19\textwidth} >{\RaggedRight\arraybackslash}X >{\RaggedRight\arraybackslash}p{0.30\textwidth}}
    \toprule
    \textbf{Tahap} & \textbf{Pertanyaan rapat evaluasi} & \textbf{Keputusan} \\
    \midrule
    Kumpulkan & Data apa yang belum masuk atau belum valid? & Lengkapi dan verifikasi \\
    Baca & Subtes, kelas, atau siswa mana yang tertinggal? & Tetapkan prioritas \\
    Tindak lanjuti & Remedial, pengayaan, atau konsultasi apa yang dibutuhkan? & Tetapkan penanggung jawab \\
    Periksa ulang & Apakah tindakan memperbaiki hasil? & Lanjutkan, ubah, atau hentikan \\
    \bottomrule
  \end{tabularx}}
\end{frame}

\begin{frame}{Orang tua menerima informasi yang relevan}
  \kicker{Pelaporan Tanpa Membandingkan Anak}
  \begin{columns}[T,onlytextwidth]
    \begin{column}{0.46\textwidth}
      {\color{StudentBlue}\bfseries Informasi utama}\par
      \begin{itemize}
        \item Kehadiran dan ketuntasan.
        \item Tren per subtes.
        \item Kekuatan dan prioritas perbaikan.
        \item Tindak lanjut yang sedang berjalan.
      \end{itemize}
    \end{column}
    \begin{column}{0.48\textwidth}
      {\color{StudentBlue}\bfseries Cara menggunakan}\par
      \begin{itemize}
        \item Menjaga rutinitas dan kesehatan.
        \item Membuka dialog tentang hambatan.
        \item Mendukung remedial tanpa menghakimi.
        \item Hadir saat keputusan program studi dibahas.
      \end{itemize}
    \end{column}
  \end{columns}
  \vspace{0.16cm}
  \claim{Tujuan laporan adalah membantu anak bergerak maju dari titik awalnya.}
\end{frame}

\begin{frame}{Prioritas perbaikan disusun bertahap}
  \kicker{Rekomendasi Hasil Evaluasi}
  {\fontsize{10.1}{12.6}\selectfont
  \renewcommand{\arraystretch}{0.98}
  \begin{tabularx}{\textwidth}{>{\bfseries\color{StudentBlue}}p{0.17\textwidth} >{\RaggedRight\arraybackslash}p{0.31\textwidth} >{\RaggedRight\arraybackslash}X}
    \toprule
    \textbf{Prioritas} & \textbf{Perbaikan} & \textbf{Hasil yang diharapkan} \\
    \midrule
    1. Data dasar & Finalisasi kohor, baseline, dan target & Denominator dan titik awal valid \\
    2. Proses & Standarkan pencatatan materi, tryout, dan konsultasi & Aktivitas dapat ditelusuri \\
    3. Tindak lanjut & Hubungkan kesalahan, remedial, dan tes ulang & Perbaikan dapat dibuktikan \\
    4. Hasil & Verifikasi penerimaan per siswa dan jalur & Persentase tidak ganda atau bias \\
    5. Komunikasi & Laporan berkala untuk sekolah, siswa, dan orang tua & Keputusan lebih cepat dan selaras \\
    \bottomrule
  \end{tabularx}}
\end{frame}

\begin{frame}{Kesimpulan evaluasi SMAN 1 Muntilan}
  \kicker{Pelajaran Utama 2026}
  \begin{columns}[T,onlytextwidth]
    \begin{column}{0.30\textwidth}
      \begin{block}{Fondasi kuat}
        Siklus materi, latihan, tryout, analisis, pembahasan, dan konsultasi sudah terbentuk.
      \end{block}
    \end{column}
    \begin{column}{0.30\textwidth}
      \begin{block}{Capaian nyata}
        Konsultasi individual dan penerimaan lintas bidang telah terdokumentasi.
      \end{block}
    \end{column}
    \begin{column}{0.30\textwidth}
      \begin{block}{Tugas 2027}
        Satukan data agar target, kemajuan, tindak lanjut, dan hasil dapat dihitung dengan benar.
      \end{block}
    \end{column}
  \end{columns}
  \vspace{0.28cm}
  \centering{\color{StudentNavy}\bfseries\fontsize{14}{17}\selectfont Program 2027 dibangun dari kekuatan 2026 dan diperbaiki melalui data yang lebih utuh.}
\end{frame}
